% page 221 of the facsimile (image p-220 of the scan)
% block 172.7 x 250.6 mm ; left margin 29.3 mm
\pgc{7.620mm}{0.000mm}{
\l{\cel{57}213}
\l{}
\l{\sou{herkogamio}. (\sou{bot.}) Nomo donita da Axell por signifikar la tipo de}
\l{\cel{3}flori qui esas organizita tale ke la autopolenizo impedesas e ke}
\l{\cel{3}la transporto di la poleno adsur la stigmato povas eventar nur}
\l{\cel{3}per la insekti.}
\l{}
\l{\sou{herkulo}. I. Mi-deo di la Panteono Romana, simbolo di forteso e vi-}
\l{\cel{3}goro. - II. (\sou{metaf.}) Persono (generale viro) di qua la forteso}
\l{\cel{3}o vigoro korpala esas tre granda. - DEFIRS.}
\l{}
\l{\sou{hermafrodita}. (Homo) a qua onu atribuas la du sexui ; (planto)}
\l{\cel{3}qua unionas en su la organi di la du sexui. - DEFIRS.}
\l{}
\l{\sou{hermetika}. Dicesas pri la klozeso di recipiento, klozeso perfekta}
\l{\cel{3}qua (ideale) obtenesas per kunfuzar la bordi di la orifico kun}
\l{\cel{3}olti di la stopilo. -DEFIRS.}
\l{}
\l{\sou{hernio.(patol.}) I. Tumoro quan produktas parto de la viceri abdo-}
\l{\cel{3}minala qua ekiras tra orifico irga, kunportante parto di la}
\l{\cel{3}peritoneo, ica\cel{2}formacante quaze sako en qua lu su lokizas.}
\l{\cel{3}- II. Tumoro quan produktas la ekiro di vicero o di parto de}
\l{\cel{3}vicero, de lua kavajo naturala tra aperturo irgequala, aciden-}
\l{\cel{3}tala o naturala. - DEFIS.}
\l{}
\l{\sou{heroo}. I. (\sou{mitol}.)Mi-deo. - II. La viro qua su distingas per sua}
\l{\cel{3}prodaji, sua granda valoro, sua devoteso. - III. La viro qua,}
\l{\cel{3}ye ta o ca instanto, akaparas la atenco, la admiro, da omna}
\l{\cel{3}kunhomi. - DEFIRS.}
\l{}
\l{\sou{herono}.-(\sou{zool.}) Granda ucelo ek la klaso \textquotedbl{}steltieri\textquotedbl{}, di qua la}
\l{\cel{3}beko esas tre longa, la gambi magra e tre longa, e qua su}
\l{\cel{3}nutras generale ek fishi. - DEFI.}
\l{}
\l{\sou{herpeto}. (patol.) Morbo konstitucionala ed heredebla, nekontagiala,}
\l{\cel{3}quan karakterizas desordini nervala. - DEFIS.}
\l{}
\l{\sou{herso}. I. (\sou{agrokultivo}).Instrumento provizita ye denti fera o ligna,}
\l{\cel{3}quan onu tranas sur sulo, pos plugo, por ruptar la glebi - pos}
\l{\cel{3}semo, por kovrar la semini. - II. (\sou{fortifikajo)}.\cel{2}Greto fera}
\l{\cel{3}o ligna, garnisita infre ye pinti grosa, quan onu elevas od}
\l{\cel{3}abasas segun-arbitrie an la pordo di fortifikajo. - F.}
\l{}
\l{\sou{heteroblastika}. (\sou{biol.})\cel{2}(Fenomeno) qua konsistas ye la formaceso,}
\l{\cel{3}en la sama punti, che animali inter-parenta, di organi inter-}
\l{\cel{3}equivalanta, ma di qui la origino embriologiala esas tamen inter-}
\l{\cel{3}diferanta.}
\l{}
\l{hetedoroxo. Qua deviacas de la ortodoxeso. - DEF.}
\l{}
\l{\sou{heterofila}. (\sou{bot.})\cel{2}(Planto) di qua la folii o la kalico-folio}
\l{\cel{3}esas kelke polimorfa.}
\l{}
\l{heterogama. (\sou{biol.}) Di qua la du elementi sexuala interdiferanta}
\l{\cel{3}esas fuzita.}
\l{}
\l{heterogena. Qua esas di naturo altra kam to a quo lu esas unio-}
\l{\cel{3}nita. - DEFIRS.}
}
